% SHARD-ID: RC-02-DEFINITIONS
% SHARD-TITLE: Definitions
% SHARD-SUMMARY: Every concept used in the book defined once, de-duplicated across the seven notes; each carries a row in db/definitions.tsv naming the notes it merges.
% SHARD-SUMMARY: The definitions are stipulated unless marked quoted, in which case a byte-checked provenance row backs them.
% SHARD-KEYWORDS: definitions, side A, side B, ring norm, Ramanujan, Ihara zeta, Hashimoto, quantum expander, Riemann channel, Bost-Connes, Weil-LPS, Artin-Schreier
% SHARD-NOTES: none
% SHARD-SCRIPTS: none

\section{Definitions}
\label{sec:definitions}

Each definition appears exactly once in the book. The column
\texttt{merged\_from} of \texttt{db/definitions.tsv} lists the notes and
sections whose wordings were merged; where the wordings differed in
substance, the difference is stated in the definition.

\subsection{The two sides}

\begin{definition}[Side A and side B]
\label{def:side-a-side-b}
For a zeta function written as $\exp\sum_{\wl\ge1}\Ncount{\wl}\uz^{\wl}/\wl$
with nonnegative counts $\Ncount{\wl}$, \emph{side A} is the gas of rings:
the counts, their Euler product over primitive closed orbits, and the
product state they form. \emph{Side B} is a transfer operator whose power
sums reproduce the counts, $\Ncount{\wl} = \Tr\,\Hash^{\wl}$ (or a
semigroup $\Tr e^{\ringlen\flowgen}$ in continuous length); its eigenvalues
are the poles and zeros.
\end{definition}

\begin{definition}[The one identity]
\label{def:one-identity}
For any matrix $M$ with eigenvalues $\lambda_1,\dots,\lambda_N$,
$1/\det(1-\uz M) = \prod_i(1-\lambda_i\uz)^{-1} =
\exp\big(\sum_{\wl\ge1}\Tr(M^{\wl})\,\uz^{\wl}/\wl\big)$. The middle term is a
bosonic partition function with one mode per eigenvector; the right term a
sum over closed walks with the $1/\wl$ of a ring's rotation symmetry.
\end{definition}

\begin{definition}[Small matrix]
\label{def:small-matrix}
For a shift or permutation with fixed-point counts
$\operatorname{fix}(\wl)$ of its $\wl$-th power, a \emph{small matrix} is any matrix
$M_{\mathrm{small}}$ with $\Tr M_{\mathrm{small}}^{\wl} =
\operatorname{fix}(\wl)$ for all $\wl$. It exists exactly when the
zeta function is rational, and its eigenvalues are not those of the
permutation matrix.
\end{definition}

\begin{definition}[RH analogue, Ramanujan property]
\label{def:rh-analogue}
A side-B operator with leading eigenvalue $\lambda_{\max}$ (or leading
growth $\qs$ per step) satisfies the \emph{RH analogue} when every other
nonzero eigenvalue has modulus $\sqrt{\lambda_{\max}}$ (respectively
$\sqrt{\qs}$). In counting terms the fluctuation of the counts is the square
root of the leading term. For graphs this is the Ramanujan property; for the
Riemann zeta it is the Riemann hypothesis; for an MPS it is one correlation
length equal to $2/\ent$.
\end{definition}

\subsection{Graphs}

\begin{definition}[Non-backtracking walk, prime cycle]
\label{def:nonbacktracking-walk}
A walk in a graph is non-backtracking if it never immediately retraces an
edge; closed in the strong sense means the last step is not the reverse of
the first. A \emph{prime cycle} $\pcycle$ is a closed non-backtracking walk that
is not a repetition of a shorter one, taken up to choice of starting point;
$\wl(\pcycle)$ is its length. On a bouquet with $\Dk/2$ loops the directed edges
are $[\Dk]$ and a closed non-backtracking walk is a cyclically reduced word.
\end{definition}

\begin{definition}[Ihara zeta]
\label{def:ihara-zeta}
For a finite graph $X$ with $\Ncount{\wl}$ closed non-backtracking walks of
length $\wl$, $\zetaX(\uz) = \exp\big(\sum_{\wl}\Ncount{\wl}\uz^{\wl}/\wl\big)
= \prod_{\pcycle}(1-\uz^{\wl(\pcycle)})^{-1}$.
\end{definition}

\begin{definition}[Hashimoto operator]
\label{def:hashimoto-operator}
The \emph{Hashimoto} (non-backtracking edge) operator of a graph is the
matrix $\Hash$ on directed edges with $\Hash_{ef} = 1$ when $f$ can follow $e$
without backtracking; then $\Ncount{\wl} = \Tr\Hash^{\wl}$ and
$\zetaX(\uz) = \det(1-\uz\Hash)^{-1}$.
\end{definition}

\begin{definition}[Ihara--Bass formula]
\label{def:ihara-bass}
For a $(\qs+1)$-regular graph with adjacency matrix $\Aadj$ on $|V|$
vertices and $|E|$ undirected edges, $\det(1-\uz\Hash) =
(1-\uz^2)^{|E|-|V|}\det(1-\uz\Aadj+\qs\uz^2)$. Each adjacency eigenvalue
$\lambda$ produces two eigenvalues $\edgeev,\edgeev'$ of $\Hash$ with $\edgeev\edgeev' = \qs$,
$\edgeev+\edgeev' = \lambda$; the eigenvalue $\qs+1$ gives $\edgeev\in\{\qs,1\}$.
\end{definition}

\begin{definition}[Ramanujan graph]
\label{def:ramanujan-graph}
A $(\qs+1)$-regular graph is \emph{Ramanujan} when every adjacency eigenvalue
other than $\pm(\qs+1)$ satisfies $|\lambda|\le2\sqrt{\qs}$, the spectral
radius of the universal covering tree. Equivalently all nontrivial $\edgeev$
satisfy $|\edgeev| = \sqrt{\qs}$, which is the RH analogue for $\zetaX$.
\end{definition}

\begin{definition}[Artin--Ihara $L$-function]
\label{def:artin-ihara}
For a graph with a representation $\hol$ of its fundamental group assigned
to closed walks (a holonomy), $\Lcurve(\uz,\hol) =
\prod_{\pcycle}\det\big(1-\hol(\pcycle)\uz^{\wl(\pcycle)}\big)^{-1}$, where the
product is over prime cycles. For a bouquet the fundamental group is free and
$\hol$ is a choice of one operator per loop.
\end{definition}

\subsection{Channels and matrix product states}

\begin{definition}[Channel, Kraus operators]
\label{def:channel}
A \emph{channel} on $\Mn$ is a completely positive trace-preserving map
$\chan(\dm) = \sum_i\Kr{i}\dm\Kr{i}^\dagger$ with $\sum_i\Kr{i}^\dagger\Kr{i} =
1$; the $\Kr{i}$ are its Kraus operators. It is \emph{unital} when
$\chan(1) = 1$ and \emph{mixed-unitary} when $\Kr{i} = \Dk^{-1/2}\Un{i}$ with
unitary $\Un{i}$, in which case $\chan(\dm) = \frac1{\Dk}\sum_i\Un{i}\dm
\Un{i}^\dagger$ and the family is taken closed under adjoint.
\end{definition}

\begin{definition}[Transfer matrix of an MPS]
\label{def:transfer-matrix}
For a matrix product state with local tensors $\Kr{1},\dots,\Kr{d}\in\Mn$
(physical dimension $d$, bond dimension $n$), the \emph{transfer matrix} is
$\Tmat = \sum_i\Ad(\Kr{i}) = \Sig$, acting on $\Mn$.
\end{definition}

\begin{definition}[Ring norm]
\label{def:ring-norm}
The \emph{ring norm} of the periodic $n$-site state $\mps{n}$ built from the
tensors is $\ringnorm{n} = \Tr\,\Tmat^{n}$; the zeta function of the state is
the generating function $\exp\sum_n\ringnorm{n}\uz^n/n = \det(1-\uz\Tmat)^{-1}$.
\end{definition}

\begin{definition}[Canonical form]
\label{def:canonical-form}
An MPS is in (left) canonical form when $\sum_k\Kr{k}^\dagger\Kr{k} = 1$,
i.e.\ when $\Tmat$ is the transfer matrix of a channel; then $\Tmat$ has
leading eigenvalue $1$ and the remaining eigenvalues lie in the closed unit
disc.
\end{definition}

\begin{definition}[Correlation length]
\label{def:correlation-length}
In canonical form write the subleading eigenvalues as $\lambda_k =
e^{-1/\corrlen{k}+i\phase{k}}$; the $\corrlen{k}$ are the correlation
lengths and the $\phase{k}$ the phases. The entropy per site is $\ent =
\log\lambda_{\max}$ of the unnormalised transfer matrix.
\end{definition}

\begin{definition}[Quantum expander]
\label{def:quantum-expander}
A channel with $\Dk$ Kraus operators is a $(n,\Dk,\lambda)$ \emph{quantum
expander} when $1$ is its unique fixed point (up to scale) and every other
eigenvalue has modulus at most $\lambda<1$.
\end{definition}

\begin{definition}[Ramanujan quantum expander]
\label{def:ramanujan-channel}
A quantum expander is \emph{Ramanujan} when $\lamtwo\le\lamH$, where, in
Hastings' words, \texttt{\textbackslash lambda\_\{H\}=\textbackslash
frac\{2\textbackslash sqrt\{D-1\}\}\{D\}.} (provenance
\texttt{prov:hastings07-bound}); in adjacency units, $\Dk\lamtwo\le2\sqrt{\Dk-1}$.
\end{definition}

\begin{definition}[Quantum Hashimoto operator, quantum Ihara zeta]
\label{def:quantum-hashimoto}
For superoperators $\Eop{1},\dots,\Eop{\Dk}$ on $\Vsp$ and a fixed-point-free
reversal $\rev$ on $[\Dk]$, the \emph{quantum Hashimoto operator} is
$\Hash$ on $\Wsp = \Vsp\otimes\mathbb C^{\Dk}$, $\Hash(v\otimes|i\rangle) =
\sum_{j\ne\bari{i}}\Eop{i}v\otimes|j\rangle$, and $\zetaT(\uz) =
\det(1-\uz\Hash)^{-1}$. For a mixed-unitary channel, $\Eop{i} = \Ad(\Un{i})$
with $\Un{\bari{i}} = \Un{i}^\dagger$, and the zeta is written $\zetaq$: the
\emph{quantum Ihara zeta} of the channel. It is the Artin--Ihara $L$-function
of the bouquet twisted by $\hol = \Ad\circ\Un{}$.
\end{definition}
